\documentclass{standalone}
\usepackage{circuitikz}
\usetikzlibrary{calc}
\definecolor{circuitnotemark}{HTML}{FE00FE}
\makeatletter
\pgfdeclareshape{smacoax}{
  \anchor{center}{\pgfpointorigin}
  \anchor{text}{\pgfpointorigin}
  \anchor{north}{\pgfpoint{0cm}{0.3cm}}
  \anchor{south}{\pgfpoint{0cm}{-0.3cm}}
  \anchor{east}{\pgfpoint{0.3cm}{0cm}}
  \anchor{west}{\pgfpoint{-0.3cm}{0cm}}
  \anchor{pin 1}{\pgfpoint{-0.7cm}{0cm}}
  \anchor{pin 2}{\pgfpoint{0cm}{-0.7cm}}
  \anchor{bpin 1}{\pgfpoint{-0.3cm}{0cm}}
  \anchor{bpin 2}{\pgfpoint{0cm}{-0.3cm}}
  \backgroundpath{
    \pgfpathmoveto{\pgfpoint{-0.274cm}{-0.122cm}}
    \pgfpatharc{204}{516}{0.3cm}
    \pgfpathmoveto{\pgfpoint{-0.7cm}{0cm}}\pgfpathlineto{\pgfpointorigin}
    \pgfpathmoveto{\pgfpoint{0cm}{-0.7cm}}\pgfpathlineto{\pgfpoint{0cm}{-0.3cm}}
  }
  \foregroundpath{
    \pgfpathcircle{\pgfpointorigin}{0.07cm}
    \pgfusepath{fill}
  }
}
\makeatother
\begin{document}
\begin{circuitikz}[american, line width=0.8pt]
\ctikzset{bipoles/length=1.2cm}
\ctikzset{grounds/scale=1.36}
\coordinate (b2) at (2,-2);
\coordinate (b3) at (4,-2);
\coordinate (b5) at (8,-2);
\coordinate (d6) at (10,-6);
\coordinate (g6) at (10,-12);
\coordinate (b6) at (10,-2);
\coordinate (b9) at (16,-2);
\coordinate (d9) at (16,-6);
\coordinate (g9) at (16,-12);
\coordinate (b10) at (18,-2);
\coordinate (b12) at (22,-2);
\coordinate (b13) at (24,-2);
\coordinate (c2) at (2,-4);
\coordinate (c13) at (24,-4);
\node[smacoax, draw, xscale=-1] (part-J1) at (b2) {}; % line 3
\node[anchor=south] at (part-J1.north) {$J_{1}$}; % line 3
\draw (b3) to[TL, l_=$\mathrm{IN}$, a^=$50\,\Omega$] (b5); % line 4
\draw (d6) to[TL, l_=$\mathrm{H1}$, a^=$73\,\Omega$] (g6); % line 5
\draw (b6) to[C, l_=$C_{\mathrm{c}}$] (b9); % line 6
\draw (d9) to[TL, l_=$\mathrm{H2}$, a^=$73\,\Omega$] (g9); % line 7
\draw (b10) to[TL, l_=$\mathrm{OUT}$, a^=$50\,\Omega$] (b12); % line 8
\node[smacoax, draw] (part-J2) at (b13) {}; % line 9
\node[anchor=south] at (part-J2.north) {$J_{2}$}; % line 9
\node[ground] at (c2) {}; % line 10
\node[ground] at (c13) {}; % line 11
\draw (part-J1.pin 1) -- (b3); % line 13
\draw (b5) -- (b6); % line 14
\draw (b6) -- (d6); % line 15
\draw (b9) -- (b10); % line 16
\draw (b9) -- (d9); % line 17
\draw (b12) -- (part-J2.pin 1); % line 18
\draw (part-J1.pin 2) -- (c2); % line 19
\draw (part-J2.pin 2) -- (c13); % line 20
\node[circ] at (b6) {};
\node[circ] at (b9) {};
\path (12,-12.593) rectangle (12.889,-11.407); % line 22
\node[anchor=west, inner sep=0, circuitnotemark, font=\large] at (12,-12) {X}; % line 22
\path (18,-12.593) rectangle (18.889,-11.407); % line 23
\node[anchor=west, inner sep=0, circuitnotemark, font=\large] at (18,-12) {X}; % line 23
\node[anchor=south west, inner sep=2pt, circuitnotemark, font=\LARGE\bfseries] (circuittitle) at (current bounding box.north west) {X};
\path (circuittitle.south west) rectangle ++(12.314,0.853);
\end{circuitikz}
\end{document}
